\section{Prompts}
\label{app:prompts}

Both prompts of \judge{NuggetRecall}, verbatim. Text in braces is filled in at run
time. Each is sent as a single user message at temperature 0.

\paragraph{Answer key.} Called once per topic, with \texttt{min\_n}=8 and
\texttt{max\_n}=20.
\begin{lstlisting}
You are building the answer key for evaluating responses to an information request.

REQUEST:
{request}

List the key facts that a complete, high-quality response to this request should contain. Each item must be:
- one self-contained factual statement, a single idea per line
- specific enough that you could check whether a response states it
- relevant to what the request actually asks for

Mark an item VITAL if a response that omits it is clearly incomplete, or OKAY if it is useful supporting detail.

Give between {min_n} and {max_n} items, most important first. Output only the list, one item per line, in exactly this format and nothing else:
1. [VITAL] <fact>
2. [OKAY] <fact>
\end{lstlisting}

\paragraph{Judging.} Called once per response. The facts are numbered from 1, one per
line.
\begin{lstlisting}
You are checking which facts a response covers. Judge only what the response says. Never use outside knowledge.

REQUEST:
{request}

RESPONSE TO JUDGE:
{answer}

FACTS TO CHECK (one per line, numbered):
{facts}

For every numbered fact, decide whether the RESPONSE conveys it:
  support  - the response clearly states this fact
  partial  - the response states part of it, gestures at it, or hedges it
  no       - the response does not convey it

Judge meaning, not wording: the same fact stated in different words is support. A fact the response merely names as a topic, without asserting it, is partial.

Output exactly one line per fact, in order, and nothing else:
1: support
2: no
3: partial

Output {count} lines, numbered 1 to {count}.
\end{lstlisting}

\paragraph{Retry suffixes.} The first attempt sends the prompt as is. The second
attempt appends the first line below and the third attempt the second, each after
a blank line. Both prompts use the same suffixes.
\begin{lstlisting}
Reminder: follow the output format exactly, with no other text.
Important: output only the numbered lines, exactly as specified, nothing else.
\end{lstlisting}
